\documentclass[11pt]{amsart}
\usepackage[letterpaper,margin=1in]{geometry}
\usepackage{amsmath,amssymb,amsthm,mathtools}
\usepackage{microtype}
\linespread{0.95}
\usepackage{enumitem}
\usepackage{needspace}
\usepackage[colorlinks=true,linkcolor=blue,citecolor=blue,urlcolor=blue]{hyperref}
\hypersetup{
  pdftitle={An Explicit Removal Estimate for Reciprocal-Sum-Free Sets},
  pdfauthor={Will Howe, Grant Huffman, Jensen Kohlmeyer, Liam Kruer},
  pdfsubject={Extremal sets of integers avoiding an equation in unit fractions},
  pdfkeywords={unit fractions, extremal sets, density, sieve methods}}

\newtheorem{theorem}{Theorem}[section]
\newtheorem{proposition}[theorem]{Proposition}
\newtheorem{lemma}[theorem]{Lemma}
\newtheorem{corollary}[theorem]{Corollary}
\theoremstyle{definition}
\newtheorem{definition}[theorem]{Definition}
\theoremstyle{remark}
\newtheorem{remark}[theorem]{Remark}
\numberwithin{equation}{section}
\newcommand{\N}{\mathbb N}
\newcommand{\Q}{\mathbb Q}
\newcommand{\R}{\mathbb R}
\newcommand{\Z}{\mathbb Z}
\newcommand{\eps}{\varepsilon}
\newcommand{\Om}{\Omega}
\newcommand{\ind}{\mathbf 1}
\newcommand{\Sp}{\mathcal S}
\newcommand{\cA}{\mathcal A}
\newcommand{\floor}[1]{\left\lfloor #1\right\rfloor}
\newcommand{\ceil}[1]{\left\lceil #1\right\rceil}
\DeclareMathOperator{\lcm}{lcm}
\allowdisplaybreaks[1]
\clubpenalty=10000
\widowpenalty=10000
\displaywidowpenalty=10000
\emergencystretch=1.5em

\title[Reciprocal-sum-free sets]{An Explicit Removal Estimate for\\Reciprocal-Sum-Free Sets}
\author[W. Howe, G. Huffman, J. Kohlmeyer, and L. Kruer]{Will Howe, Grant Huffman, Jensen Kohlmeyer, and Liam Kruer}
\date{November 11, 2026}
\thanks{\textit{E-mail addresses}: \href{mailto:howew@purdue.edu}{\texttt{howew@purdue.edu}} (W. Howe),
\href{mailto:huffma97@purdue.edu}{\texttt{huffma97@purdue.edu}} (G. Huffman),
\href{mailto:jenwin@purdue.edu}{\texttt{jenwin@purdue.edu}} (J. Kohlmeyer),
\href{mailto:kruerl@purdue.edu}{\texttt{kruerl@purdue.edu}} (L. Kruer).}
\subjclass[2020]{11D68, 11N35}
\keywords{unit fractions, extremal sets, density, sieve methods, effective approximation}

\begin{document}

\begin{abstract}
Let $f(N)$ be the largest cardinality of a subset of $\{1,\ldots,N\}$ containing no three distinct elements with $1/a=1/b+1/c$.
We give an explicit uniform removal estimate for finite-prime relaxations of this problem.
Together with the smooth-number decomposition, the estimate yields $f(N)=CN+o(N)$ and an effective description of $C$ by finite optimization problems.
If $U_V$ is the limiting density obtained by retaining triples with smooth primitive coordinates over the primes at most $2V$, then
\[
0\le U_V-C\le (10^{23}+561000)(\log V)^{-1/140}\qquad(V\ge4).
\]
The proof uses elementary prime estimates, a finite quadratic sieve, and a dyadic decomposition indexed by the larger primitive parameter.
We also give rational enclosures of arbitrary prescribed width and an explicit modulus of convergence for $f(N)/N$.
The finite-prime framework and the underlying sieve are due to the earlier density-limit argument; the emphasis here is the uniform quantitative estimate.
\end{abstract}

\maketitle

\section{Introduction and main results}

A finite set $\cA\subset\N$ is \emph{reciprocal-sum-free} if it contains no pairwise distinct $a,b,c$ satisfying
\begin{equation}\label{eq:reciprocal}
\frac1a=\frac1b+\frac1c.
\end{equation}
Throughout, $\N=\{1,2,\ldots\}$, and
\[
f(N)=\max\bigl\{|\cA|:\cA\subseteq\{1,\ldots,N\},\ \cA\text{ is reciprocal-sum-free}\bigr\}.
\]
Estimating $f(N)$ is Erd\H{o}s problem 302, originating in the unit-fraction questions of Erd\H{o}s and Graham \cite{EG}.
The existence of its limiting density and the finite-prime description used below occur in the September 2026 manuscript accompanying \cite{CGM}.
We develop that framework with explicit constants and a removal estimate uniform over the whole interval $\{1,\ldots,N\}$.

For a nonempty finite set $P$ of primes, write
\[
\Sp(P)=\{n\in\N:p\mid n\Longrightarrow p\in P\},
\qquad
\rho_P=\prod_{p\in P}\left(1-\frac1p\right).
\]
The integer $1$ belongs to $\Sp(P)$.
Let $\alpha_P(t)$ be the largest size of a reciprocal-sum-free subset of $\Sp(P)\cap[1,t]$, with $\alpha_P(0)=0$.
For $V\ge4$, put
\begin{equation}\label{eq:Pdef}
P_V=\{p\text{ prime}:p\le2V\}.
\end{equation}
If $1=s_1<s_2<\cdots$ are the elements of $\Sp(P_V)$, define
\begin{equation}\label{eq:Udef}
U_V=\rho_{P_V}\sum_{j\ge1}\frac{\alpha_{P_V}(s_j)-\alpha_{P_V}(s_{j-1})}{s_j},
\qquad \alpha_{P_V}(s_0)=0.
\end{equation}
The series converges absolutely, as proved in Section~\ref{sec:fibres}.

\begin{theorem}\label{thm:density}
The sequence $U_V$ is nonincreasing and
\begin{equation}\label{eq:Cdef}
\lim_{N\to\infty}\frac{f(N)}N=C,
\qquad
C=\inf_{V\ge4}U_V.
\end{equation}
The constant $C$ admits rational enclosures of every positive prescribed width. Moreover, $C\ge5/8$.
\end{theorem}

The quantitative result concerns a relaxation of the forbidden equation.
For a triple $(a,b,c)$ of positive integers, its \emph{primitive coordinates} are the three entries after division by $\gcd(a,b,c)$.
Let $F_V(N)$ be the largest size of a subset of $\{1,\ldots,N\}$ containing no distinct solution of \eqref{eq:reciprocal} whose primitive coordinates belong to $\Sp(P_V)$.
Thus $f(N)\le F_V(N)$.

\begin{theorem}[Uniform removal estimate]\label{thm:removal}
Set $z=561/500$ and $q=20/13$.
For all integers $N,V\ge4$ and $K\ge0$,
\begin{equation}\label{eq:removal}
\begin{split}
0\le F_V(N)-f(N)
\le N\bigl[&500000z^{-K}\\
&+10^{23}q^{3K}\bigl((\log V)^{-7/80}+(\log N)^{-7/80}\bigr)\bigr].
\end{split}
\end{equation}
\end{theorem}

\begin{corollary}\label{cor:rates}
For every integer $V\ge4$,
\begin{equation}\label{eq:outerrate}
0\le U_V-C\le (10^{23}+561000)(\log V)^{-1/140}.
\end{equation}
For integers $N,V\ge4$,
\begin{equation}\label{eq:finiterate}
0\le\frac{F_V(N)-f(N)}N
\le(2\cdot10^{23}+561000)\bigl(\log\min(V,N)\bigr)^{-1/140}.
\end{equation}
\end{corollary}

All logarithms are natural.
The constants are chosen to make every estimate uniform, including the smallest admissible cutoffs.
The finite-prime formula and the quadratic sieve in Section~\ref{sec:sieve} follow the earlier argument in \cite{CGM}.
The organization of the remaining triples by their larger primitive parameter permits a single regularity condition for all primes.
Prime-factor regularity also appears in the unit-fraction work of Della Pietra \cite{DP}.
The argument below includes proofs of the estimates it uses.

\section{Primitive triples and finite-prime fibres}\label{sec:fibres}

\begin{lemma}\label{lem:parameters}
Every positive solution of \eqref{eq:reciprocal} with $b<c$ has a unique representation
\begin{equation}\label{eq:parameters}
(a,b,c)=\bigl(guv,\ gu(u+v),\ gv(u+v)\bigr),
\quad g\ge1,\quad 1\le u<v,\quad \gcd(u,v)=1.
\end{equation}
The three coordinates inside the common factor $g$ have greatest common divisor $1$.
Conversely, every tuple in \eqref{eq:parameters} is a solution with $a<b<c$.
\end{lemma}

\begin{proof}
Write $b=hu$ and $c=hv$, where $h=\gcd(b,c)$ and $\gcd(u,v)=1$.
Clearing denominators gives $a=huv/(u+v)$.
Since $\gcd(uv,u+v)=1$, the integer $u+v$ divides $h$.
Writing $h=g(u+v)$ proves the representation.
The integers $u,v,u+v$ are pairwise coprime, so the three pairwise products have greatest common divisor $1$.
The reduced ratio $b/c=u/v$ determines $u,v$, after which $g$ is determined.
The converse follows by substitution, and the inequalities follow from $0<u<v$.
\end{proof}

An integer is \emph{$P$-rough} if it is divisible by no prime in $P$.
Unique factorization gives a unique expression $n=rs$, with $r$ $P$-rough and $s\in\Sp(P)$.
Let $F_P(N)$ denote the analogous relaxation for a nonempty finite prime set $P$, and let
\[
R_P(x)=\#\{r\in\N:r\le x,\ r\text{ is }P\text{-rough}\}.
\]

\Needspace{11\baselineskip}
\begin{lemma}\label{lem:fibreidentity}
Write $1=s_1<s_2<\cdots$ for $\Sp(P)$, and set
$\Delta_{P,j}=\alpha_P(s_j)-\alpha_P(s_{j-1})$, with the same initial convention as above.
Then $\Delta_{P,j}\in\{0,1\}$ and
\begin{equation}\label{eq:fibreidentity}
F_P(N)=\sum_{j\ge1}\Delta_{P,j}R_P(N/s_j).
\end{equation}
Consequently,
\begin{equation}\label{eq:fixedlimit}
\lim_{N\to\infty}\frac{F_P(N)}N
=\rho_P\sum_{j\ge1}\frac{\Delta_{P,j}}{s_j}.
\end{equation}
\end{lemma}

\begin{proof}
Adding one point to a finite set increases its maximum reciprocal-sum-free cardinality by either zero or one.
This proves the assertion about $\Delta_{P,j}$.

A retained triple has $P$-smooth primitive coordinates, so its three entries have the same $P$-rough part.
Conversely, if a reciprocal triple lies in one fibre $r\Sp(P)$, division by its common divisor leaves $P$-smooth coordinates.
The forbidden triples in the fibre $r\Sp(P)\cap[1,N]$ therefore correspond exactly to the original forbidden triples in $\Sp(P)\cap[1,N/r]$.
Different fibres impose independent conditions. Hence
\[
F_P(N)=\sum_{\substack{r\le N\\r\ P\text{-rough}}}\alpha_P(N/r).
\]
Expand $\alpha_P(t)=\sum_{s_j\le t}\Delta_{P,j}$ and interchange these finite sums to obtain \eqref{eq:fibreidentity}.

If $d=|P|$, finite inclusion--exclusion gives
\begin{equation}\label{eq:roughcount}
|R_P(x)-\rho_Px|\le2^d\qquad(x\ge0).
\end{equation}
Also $R_P(N/s_j)/N\le1/s_j$, and
\begin{equation}\label{eq:euler}
\sum_{s\in\Sp(P)}\frac1s=\prod_{p\in P}(1-p^{-1})^{-1}=\rho_P^{-1}.
\end{equation}
The identity follows by multiplying finitely many convergent geometric series.
Dominated convergence in \eqref{eq:fibreidentity} proves \eqref{eq:fixedlimit}.
\end{proof}

It follows that $F_V(N)/N\to U_V$.
Increasing $V$ adds forbidden triples, so $F_V(N)$ and $U_V$ are nonincreasing in $V$.
In particular, the infimum in \eqref{eq:Cdef} exists in $[0,1]$.

For a finite cutoff $T$, define
\begin{equation}\label{eq:prefix}
L_P(T)=\rho_P\sum_{s_j\le T}\frac{\Delta_{P,j}}{s_j},
\qquad
E_P(T)=1-\rho_P\sum_{\substack{s\le T\\s\in\Sp(P)}}\frac1s.
\end{equation}
Both quantities are rational, $E_P(T)\ge0$, and \eqref{eq:euler} gives
\begin{equation}\label{eq:prefixerror}
0\le \rho_P\sum_{j\ge1}\frac{\Delta_{P,j}}{s_j}-L_P(T)\le E_P(T),
\qquad E_P(T)\longrightarrow0.
\end{equation}

\section{Prime estimates and regular integers}\label{sec:regular}

For $n\ge1$ and $x\ge2$, let
\[
\Om(n;x)=\sum_{p\le x}v_p(n),\qquad \Om(n)=\sum_pv_p(n).
\]
We first record the elementary estimates needed in the proof.

\begin{lemma}\label{lem:primebounds}
For $x\ge2$,
\begin{equation}\label{eq:primebounds}
\log\log x-1\le\sum_{p\le x}\frac1p\le\log\log x+14.
\end{equation}
Moreover, writing $\psi(x)=\sum_{p^a\le x}\log p$, one has $\psi(x)\le4x$ for $x\ge1$ and
\begin{equation}\label{eq:primelog}
\sum_{p<y}\frac{\log p}{p}\le8\log y\qquad(y\ge2).
\end{equation}
\end{lemma}

\begin{proof}
For an integer $n\ge2$, every prime power in $(\ceil{n/2},n]$ contributes to the valuation of $\binom n{\floor{n/2}}$.
The other valuation terms are nonnegative. Thus
\[
\psi(n)-\psi(\ceil{n/2})\le\log\binom n{\floor{n/2}}\le n\log2.
\]
Induction proves $\psi(n)\le4n$, since $4\floor{n/2}\ge n\log2$ for $n\ge2$.
The real-variable assertion follows by taking integer parts.
With $\vartheta(t)=\sum_{p\le t}\log p\le4t$, partial summation from $2$ gives
\[
\sum_{p\le y}\frac{\log p}{p}
\le4+4\log(y/2)\le8\log y,
\]
where $\log2>1/2$ proves the last inequality.

The Euler product over $p\le x$ contains the harmonic sum through $\floor x$, which is at least $\log x$.
Its logarithm differs from $\sum_{p\le x}1/p$ by at most
\[
\sum_{p\le x}\sum_{a\ge2}\frac1{ap^a}
\le\sum_{n\ge2}\frac1{n(n-1)}=1.
\]
This proves the lower bound in \eqref{eq:primebounds}.

For the upper bound put $s=1+1/\log x$.
Since $1-e^{-t}\le t$ for $t\ge0$,
\[
\sum_{p\le x}\frac1p
\le\sum_p p^{-s}+(s-1)\sum_{p\le x}\frac{\log p}{p}.
\]
Partial summation with $\psi(t)\le4t$ bounds the last prime sum by $4+4\log x$.
The Euler product and the integral test give
$\sum_p p^{-s}\le\log\zeta(s)\le\log(1+\log x)$.
Finally,
\[
\log(1+\log x)+4+\frac4{\log x}
\le\log\log x+14,
\]
using $\log x>1/2$ and $\log(1+1/\log x)<2$.
\end{proof}

Fix the parameters
\begin{equation}\label{eq:params}
\alpha=\frac{53}{50},\quad z=\frac{561}{500},\quad q=\frac{20}{13},
\quad\theta=q^{-2}=\frac{169}{400},\quad\gamma=q^{-3}=\frac{2197}{8000}.
\end{equation}
An integer $n$ is \emph{$K$-regular} if
\begin{equation}\label{eq:regular}
\Om(n;x)\le\alpha\log\frac{\log x}{\log2}+K\qquad(x\ge2).
\end{equation}
It is equivalent to require this inequality only at integer $x\ge2$.

\begin{lemma}[Exceptional integers]\label{lem:exceptional}
For every $N\ge1$ and integer $K\ge0$, at most $500000z^{-K}N$ integers in $[1,N]$ fail \eqref{eq:regular}.
\end{lemma}

\begin{proof}
For $1<z<2$, the divisor expansion of $z^{\Om(n;x)}$ has nonnegative multiplicative coefficients
$h(p^a)=z^a-z^{a-1}$ for $p\le x$, and $h(p^a)=0$ otherwise.
Consequently,
\begin{equation}\label{eq:momentproduct}
\sum_{n\le N}z^{\Om(n;x)}
\le N\sum_{d\ge1}\frac{h(d)}d
=N\prod_{p\le x}\left(1+\frac{z-1}{p-z}\right).
\end{equation}
The logarithm of the product is at most
\[
(z-1)\sum_{p\le x}\frac1p
+\frac{(z-1)z}{1-z/2}\sum_p\frac1{p^2}.
\]
Use \eqref{eq:primebounds} and $\sum_p p^{-2}\le\sum_{n\ge2}[n(n-1)]^{-1}=1$.
For the specified $z$,
\[
14(z-1)+\frac{(z-1)z}{1-z/2}<\frac{56}{27}<3\log2.
\]
The last strict inequality follows from the first two terms of
$\log2=2\sum_{m\ge0}(1/3)^{2m+1}/(2m+1)$.
As $\log2<1$, we obtain
\begin{equation}\label{eq:moment}
\sum_{n\le N}z^{\Om(n;x)}
\le8N\left(\frac{\log x}{\log2}\right)^{z-1}.
\end{equation}

Put $d=\alpha\log z-z+1$.
If \eqref{eq:regular} fails at a point with
$e^j\le\log x/\log2<e^{j+1}$, it follows that
\[
\Om(n;x_j)>\alpha j+K,
\qquad x_j=\floor{\exp((\log2)e^{j+1})}.
\]
By \eqref{eq:moment} and Markov's inequality, the number of such $n$ is at most
$8Ne^{z-1}z^{-K}e^{-dj}$.
Summation over $j\ge0$ gives the bound
\[
\frac{8e^{z-1}}{1-e^{-d}}z^{-K}N.
\]
The two-term lower estimate
\[
\log z\ge2\left(\frac{61}{1061}+\frac13\left(\frac{61}{1061}\right)^3\right)
\]
implies $d>1/52000$ by rational arithmetic.
Also $e^{z-1}<500/439$, while $e^d\ge1+d$ gives
$(1-e^{-d})^{-1}\le1+d^{-1}<52001$.
Thus the coefficient is less than $208004000/439<500000$.
\end{proof}

\begin{lemma}[Common-scale mean]\label{lem:mean}
If $X\ge2$, $Y\ge1$, and $0<t\le1$, then
\begin{equation}\label{eq:mean}
\sum_{g\le Y}t^{\Om(g;X)}
\le4\cdot10^9Y\bigl((\log(2X))^{t-1}+(\log(2Y))^{t-1}\bigr).
\end{equation}
\end{lemma}

\begin{proof}
Let $G(n)=t^{\Om(n;X)}$ and $H(Y)=\sum_{n\le Y}G(n)/n$.
The function $G$ is completely multiplicative and $0<G(n)\le1$.
Writing $\log n=\sum_{d\mid n}\Lambda(d)$ and then $n=dm$, we obtain
\[
\sum_{n\le Y}G(n)\log n
=\sum_{m\le Y}G(m)\sum_{d\le Y/m}G(d)\Lambda(d)
\le4YH(Y).
\]
For $Y\ge2$, split the sum of $G(n)$ at $\sqrt Y$.
Since $\log(2Y)\le2\log Y$, $H(Y)\ge1$, and $\log(2Y)\le2\sqrt Y$, this yields
\begin{equation}\label{eq:meansplit}
\sum_{n\le Y}G(n)
\le\sqrt Y+\frac{8Y}{\log Y}H(Y)
\le\frac{18Y}{\log(2Y)}H(Y).
\end{equation}
Here $\log(2Y)\le2\sqrt Y$ follows by minimizing $2\sqrt Y-\log(2Y)$ on $[1,\infty)$.

Set $m=\min(X,Y)$.
Multiplication of geometric series bounds $H(Y)$ by
\[
\prod_{p\le Y}(1-p^{-1})^{-1}
\prod_{p\le m}\frac{1-p^{-1}}{1-tp^{-1}}.
\]
The first product is at most $e^{15}\log(2Y)$, by \eqref{eq:primebounds} and the higher-power estimate in its proof.
For the second product, the power series for the logarithm gives
\[
\log\frac{1-p^{-1}}{1-tp^{-1}}
=-\sum_{a\ge1}\frac{1-t^a}{ap^a}\le-\frac{1-t}{p}.
\]
Hence that product is at most $e(\log m)^{t-1}$.
Since $m\ge2$ and $\log(2m)\le2\log m$,
\[
(\log m)^{t-1}\le2(\log(2m))^{t-1}.
\]
Combining this with \eqref{eq:meansplit} gives
\[
\sum_{n\le Y}G(n)
\le36e^{16}Y(\log(2m))^{t-1}.
\]
The negative power at $m$ is at most the sum of the two powers in \eqref{eq:mean}, and
$36e^{16}<36\cdot3^{16}<2\cdot10^9$.
For $1\le Y<2$, the left side is $1$, while $(\log(2Y))^{t-1}\ge1/2$.
This also proves the stated bound.
\end{proof}

\section{A finite quadratic sieve}\label{sec:sieve}

The next lemma is the uniform finite sieve underlying the weighted argument in \cite[Appendix B]{CGM}.
We include the finite-dimensional proof, including the error term.

\begin{lemma}\label{lem:sieve}
Let $X>1$, $Z=X^{1/768}\ge2$, and $M_0=\floor{4X}$.
Let $\mathcal P$ be any finite set of primes less than $Z$.
For each $p\in\mathcal P$, let $U_p$ be a proper subset of $(\Z/p\Z)^2$ with $|U_p|\le3p$, and put $\eta_p=|U_p|/p^2$.
Then
\begin{equation}\label{eq:sieve}
\#\{(u,v)\in[1,M_0]^2:(u,v)\bmod p\notin U_p\text{ for all }p\in\mathcal P\}
\le32X^2\prod_{p\in\mathcal P}(1-\eta_p)+9X^{3/2}.
\end{equation}
\end{lemma}

\begin{proof}
Primes with $U_p=\varnothing$ may be removed.
Let $P_0$ be the product of the remaining primes and set $D=X^{1/16}$.
All divisors of $P_0$ are squarefree.
For such divisors define the multiplicative functions
\[
\eta(d)=\prod_{p\mid d}\eta_p,
\qquad H(d)=\prod_{p\mid d}\frac{\eta_p}{1-\eta_p},
\qquad K(d)=H(d)^{-1}.
\]
The empty products are $1$.
Write
\[
\mathcal D=\{d\mid P_0:d\le D\},\qquad
\mathcal H=\sum_{d\mid P_0}H(d)=\prod_{p\mid P_0}(1-\eta_p)^{-1},\qquad
G=\sum_{d\in\mathcal D}H(d).
\]
Under the probability measure $H(d)/\mathcal H$ on the divisors of $P_0$, each prime $p$ occurs with probability $\eta_p$.
By \eqref{eq:primelog},
\[
\frac1{\mathcal H}\sum_{d\mid P_0}H(d)\log d
=\sum_{p\mid P_0}\eta_p\log p
\le24\log Z=\frac12\log D.
\]
Markov's inequality therefore gives
\begin{equation}\label{eq:halfmass}
G\ge\mathcal H/2.
\end{equation}

For $\ell\in\mathcal D$, put $Y_\ell=\mu(\ell)H(\ell)/G$, and for $d\in\mathcal D$ define
\begin{equation}\label{eq:lambdadef}
\lambda_d=\frac1{\eta(d)}\sum_{\substack{\ell\in\mathcal D\\d\mid\ell}}\mu(\ell/d)Y_\ell.
\end{equation}
Finite divisor inversion gives
\begin{equation}\label{eq:lambdaidentity}
\sum_{\substack{d\in\mathcal D\\\ell\mid d}}\lambda_d\eta(d)=Y_\ell,
\qquad \lambda_1=1.
\end{equation}
Indeed, the coefficient after substitution is a divisor sum of the M\"obius function, equal to $1$ at the identity and $0$ otherwise.
The identities
\[
\eta(\lcm(d,e))=\frac{\eta(d)\eta(e)}{\eta(\gcd(d,e))},
\qquad \frac1{\eta(m)}=\sum_{\ell\mid m}K(\ell)
\]
now diagonalize the quadratic form:
\begin{equation}\label{eq:quadratic}
\begin{split}
\sum_{d,e\in\mathcal D}\lambda_d\lambda_e\eta(\lcm(d,e))
&=\sum_{\ell\in\mathcal D}K(\ell)
\left(\sum_{\substack{d\in\mathcal D\\\ell\mid d}}\lambda_d\eta(d)\right)^2\\
&=\sum_{\ell\in\mathcal D}K(\ell)Y_\ell^2=\frac1G.
\end{split}
\end{equation}

Writing $\ell=dm$ in \eqref{eq:lambdadef}, squarefreeness gives
$\mu(m)\mu(dm)=\mu(d)$.
Since the resulting sum of $H(m)$ is at most $G$,
\begin{equation}\label{eq:lambdabound}
|\lambda_d|\le\frac{H(d)}{\eta(d)}
=\prod_{p\mid d}(1-\eta_p)^{-1}\le d^2.
\end{equation}
The last inequality uses the integer inequality $p^2-|U_p|\ge1$.

Let $I_d(u,v)$ indicate membership in $U_p$ at every $p\mid d$, with $I_1=1$.
There are $d^2\eta(d)$ such residue pairs modulo $d$, by the Chinese remainder theorem.
Each one-dimensional residue class has $M_0/d$ representatives with error at most $1$.
Thus
\begin{equation}\label{eq:crtcount}
\sum_{1\le u,v\le M_0}I_d(u,v)=M_0^2\eta(d)+R_d,
\qquad |R_d|\le2M_0d+d^2.
\end{equation}
For a pair avoiding every $U_p$, the sum $\sum_{d\in\mathcal D}\lambda_dI_d(u,v)$ equals $1$.
Its square is nonnegative for every other pair.
Expanding the sum of squares and using \eqref{eq:quadratic}--\eqref{eq:crtcount} gives an upper bound
\[
\frac{M_0^2}{G}
+\sum_{d,e\in\mathcal D}\lambda_d\lambda_eR_{\lcm(d,e)}.
\]
There are at most $D^2$ ordered pairs $(d,e)$, their least common multiples are at most $D^2$, and \eqref{eq:lambdabound} bounds $|\lambda_d\lambda_e|$ by $D^4$.
The absolute error is therefore at most
\[
D^6(8XD^2+D^4)=8X^{3/2}+X^{5/8}\le9X^{3/2}.
\]
Finally, \eqref{eq:halfmass} gives
$M_0^2/G\le32X^2\prod_{p\in\mathcal P}(1-\eta_p)$.
The argument also applies when $\mathcal P$ is empty.
\end{proof}

\section{A weighted square estimate}\label{sec:square}

Set
\begin{equation}\label{eq:beta}
\beta=3(1-\theta)=\frac{693}{400},
\qquad d_\theta=2-3\theta+\theta^3=\frac{51706809}{64000000}.
\end{equation}

\begin{lemma}\label{lem:square}
For every integer $X\ge2$,
\begin{equation}\label{eq:square}
\sum_{1\le u,v\le2X}\theta^{\Om(u)+\Om(v)+\Om(u+v)}
\le20\cdot10^9X^2(\log(2X))^{-\beta}.
\end{equation}
\end{lemma}

\begin{proof}
First suppose $Z=X^{1/768}\ge16$, and let
\[
\mathcal P=\{p\text{ prime}:3\le p\le\floor{Z/2}\}.
\]
At every $p\in\mathcal P$, select each of the three residue lines
\[
u=0,\qquad v=0,\qquad u+v=0
\]
independently with probability $1-\theta$.
Let $U_p$ be their union.
As $p\ge3$, every realization satisfies $|U_p|\le3p-2<p^2$.
For a fixed integer pair $(u,v)$, its probability of surviving all the selections is
\[
\prod_{p\in\mathcal P}
\theta^{\ind_{p\mid u}+\ind_{p\mid v}+\ind_{p\mid u+v}}
\ge\theta^{\Om(u)+\Om(v)+\Om(u+v)}.
\]
The three lines meet only at the origin. The average surviving proportion in $(\Z/p\Z)^2$ is consequently
\begin{equation}\label{eq:localdensity}
m_p=\frac{p^2-3p+2+3(p-1)\theta+\theta^3}{p^2}
=1-\frac\beta p+\frac{d_\theta}{p^2}.
\end{equation}
Taking expectations in Lemma~\ref{lem:sieve}, independently at each prime, yields
\begin{equation}\label{eq:weightsieve}
\sum_{1\le u,v\le\floor{4X}}\theta^{\Om(u)+\Om(v)+\Om(u+v)}
\le32X^2\prod_{p\in\mathcal P}m_p+9X^{3/2}.
\end{equation}

Put $T=\floor{Z/2}$.
For $Z\ge16$, $T\ge Z/4$ and $\log T\ge\tfrac12\log Z$.
Removing the prime $2$ from the lower estimate in \eqref{eq:primebounds} gives
\[
\sum_{p\in\mathcal P}\frac1p
\ge\log\log T-\frac32
\ge\log\log Z-\frac52.
\]
The factors $m_p$ are positive, and $\log(1+t)\le t$ gives
\[
\log\prod_{p\in\mathcal P}m_p
\le-\beta\log\log Z+\frac52\beta+d_\theta.
\]
Here
\[
\frac52\beta+d_\theta=\frac{328906809}{64000000}<\frac{16}{3}\le\log256.
\]
Thus
\begin{equation}\label{eq:productdecay}
\prod_{p\in\mathcal P}m_p\le256(\log Z)^{-\beta}.
\end{equation}

For $X\ge1$, the elementary inequality
\begin{equation}\label{eq:logsqrt}
(\log X)^2\le16\sqrt X
\end{equation}
follows from $\log X\le4\exp((\log X)/4)$.
Since $\beta<2$ and $\log X\ge1$ in the present range,
$9X^{3/2}\le144X^2(\log X)^{-\beta}$.
Combining \eqref{eq:weightsieve} and \eqref{eq:productdecay}, and using $\log Z=(\log X)/768$, gives
\[
\sum_{1\le u,v\le\floor{4X}}\theta^{\Om(u)+\Om(v)+\Om(u+v)}
\le(32\cdot256\cdot768^2+144)X^2(\log X)^{-\beta}
<5\cdot10^9X^2(\log X)^{-\beta}.
\]
Restriction to $u,v\le2X$ and $\log(2X)\le2\log X$ prove \eqref{eq:square}.

It remains to treat $2\le X<16^{768}$.
There are $4X^2$ pairs and each weight is at most $1$.
Moreover, $1\le\log(2X)<3073$ and $\beta<2$.
The numerical inequality $4\cdot3073^2<20\cdot10^9$ therefore proves \eqref{eq:square} in this range as well.
\end{proof}

\section{Counting regular triples}\label{sec:count}

Set
\begin{equation}\label{eq:MH}
M=3\alpha\log q,\qquad H=M-\beta,\qquad\kappa=\frac7{80}.
\end{equation}
The following inequalities determine the available summation range:
\begin{equation}\label{eq:exponents}
0<M<2,\qquad -\frac12<H<0,\qquad H+\gamma<-\kappa.
\end{equation}
For completeness, the series
$\log q=2\sum_{m\ge0}(7/33)^{2m+1}/(2m+1)$ gives
\[
\frac{14}{33}\le\log q<\frac{4308}{10000}.
\]
For the upper bound, keep $m=0,1,2,3$ and bound the remaining terms by
$2(7/33)^9/[9(1-(7/33)^2)]$.
Substitution into \eqref{eq:MH} proves every inequality in \eqref{eq:exponents} by rational arithmetic.

\begin{lemma}[One dyadic block]\label{lem:block}
Let $X\ge2$ be an integer and $X^2\le N$.
The number of parameter triples in \eqref{eq:parameters} with
$X\le v<2X$, $gv(u+v)\le N$, and all three coordinates $K$-regular is at most
\begin{equation}\label{eq:block}
72\cdot10^{19}q^{3K}N
\left[R^{H+\gamma-1}+R^H\left(\log\frac{2N}{X^2}\right)^{\gamma-1}\right],
\qquad R=\log(2X).
\end{equation}
\end{lemma}

\begin{proof}
Write $w=u+v$ and $Y=N/X^2$.
Then $X\le w<4X$ and $g\le Y$.
Add the three inequalities \eqref{eq:regular} at cutoff $w$.
Complete additivity of $\Om(\cdot;w)$ gives
\[
3\Om(g;w)+2\bigl(\Om(u)+\Om(v)+\Om(w)\bigr)
\le3\alpha\log\frac{\log w}{\log2}+3K.
\]
Since $\gamma=q^{-3}$ and $\theta=q^{-2}$, exponentiation with base $q$ bounds the indicator that all three coordinates are regular by
\begin{equation}\label{eq:pointweight}
q^{3K}\left(\frac{\log w}{\log2}\right)^M
\gamma^{\Om(g;w)}\theta^{\Om(u)+\Om(v)+\Om(w)}.
\end{equation}
Because $\gamma<1$ and $w\ge X$, replacing $\Om(g;w)$ by $\Om(g;X)$ increases this expression.
Also $\log w/\log2\le3R$, and $0<M<2$ gives $(3R)^M\le9R^M$.

We may now sum over all $g\le Y$ and all $1\le u,v\le2X$, discarding coprimality and the remaining restrictions.
Apply Lemma~\ref{lem:mean} with $t=\gamma$, and Lemma~\ref{lem:square} to the pair $(u,v)$.
The resulting bound is
\[
9q^{3K}R^M\cdot4\cdot10^9Y
\bigl(R^{\gamma-1}+(\log(2Y))^{\gamma-1}\bigr)
\cdot20\cdot10^9X^2R^{-\beta}.
\]
Since $YX^2=N$ and $H=M-\beta$, this is \eqref{eq:block}.
\end{proof}

\begin{lemma}[Dyadic tails]\label{lem:tails}
Let $N,V\ge4$, and let $\mathcal J$ be any set of integers $j\ge1$ for which
$V<2^{j+1}$ and $2^{2j}\le N$.
With $R_j=(j+1)\log2$,
\begin{align}
\sum_{j\in\mathcal J}R_j^{H+\gamma-1}
&\le128(\log V)^{-\kappa},\label{eq:tailone}\\
\sum_{j\in\mathcal J}R_j^H
\bigl(\log N+(1-2j)\log2\bigr)^{\gamma-1}
&\le64(\log N)^{-\kappa}.\label{eq:tailtwo}
\end{align}
\end{lemma}

\begin{proof}
Put $h=\log2$, so $1/2<h<1$.
For the first estimate, write $W=\log V>1$ and $J=\floor W$.
Since $R_j>W$ and $R_j\le j+1$, every $j\in\mathcal J$ satisfies $j\ge J$.
Also $J\ge W/2$.
By \eqref{eq:exponents}, $R_j\ge1$ implies
\[
R_j^{H+\gamma-1}\le R_j^{-1-\kappa}
\le4(j+1)^{-1-\kappa}.
\]
The integral test now gives
\[
\sum_{j\in\mathcal J}R_j^{H+\gamma-1}
\le\frac4\kappa J^{-\kappa}
\le\frac8\kappa W^{-\kappa}
<128W^{-\kappa}.
\]

For the second estimate put $U=\log N>1$ and split at $jh=U/4$.
In the lower part,
\[
U+(1-2j)h\ge U/2,\qquad R_j^H\le2(j+1)^H.
\]
If this part is nonempty, $j\ge1$ forces $U\ge4h>2$, and $j+1\le U$.
Since $H>-1/2$,
\[
\sum_{\substack{j\in\mathcal J\\jh\le U/4}}(j+1)^H
\le\sum_{1\le n\le\floor U}n^H
\le\frac{U^{H+1}}{H+1}\le2U^{H+1}.
\]
As $-1<\gamma-1<0$, the lower contribution is at most $8U^{H+\gamma}$.

For the upper part let $B=\floor{U/(2h)}$.
Every permitted $j$ is at most $B$, and $B\le U$.
Here $R_j>U/4$, so $R_j^H\le4U^H$.
Furthermore,
\[
U+(1-2j)h\ge(B-j+1)h,
\]
and hence its $(\gamma-1)$st power is at most $2(B-j+1)^{\gamma-1}$.
The substitution $m=B-j+1$ is injective and takes values in $[1,B]$.
Therefore the upper contribution is at most
\[
8U^H\sum_{m=1}^B m^{\gamma-1}
\le\frac8\gamma U^{H+\gamma}
\le32U^{H+\gamma},
\]
using $\gamma>1/4$.
The two contributions total at most $40U^{H+\gamma}\le64U^{-\kappa}$, by \eqref{eq:exponents}.
The same bounds hold when either part is empty.
\end{proof}

\begin{proposition}\label{prop:regularcount}
For integers $N,V\ge4$ and $K\ge0$, the number of ordered reciprocal triples in $[1,N]$ with primitive parameter $v>V$ and all three coordinates $K$-regular is at most
\begin{equation}\label{eq:regularcount}
10^{23}q^{3K}N\bigl((\log V)^{-\kappa}+(\log N)^{-\kappa}\bigr).
\end{equation}
\end{proposition}

\begin{proof}
For each triple choose the unique $j$ with $2^j\le v<2^{j+1}$ and use Lemma~\ref{lem:block} with $X=2^j$.
As $v>V\ge4$, one has $j\ge1$ and $V<2^{j+1}$.
Also $gv(u+v)\le N$ implies $2^{2j}\le N$.
The logarithm in the second term of \eqref{eq:block} is exactly
$\log N+(1-2j)\log2$.
Apply Lemma~\ref{lem:tails} and use
\[
72\cdot10^{19}\cdot128<10^{23},
\qquad 72\cdot10^{19}\cdot64<10^{23}.
\]
Uniqueness in Lemma~\ref{lem:parameters} ensures that this bounds the number of triples.
\end{proof}

\section{Removal, convergence, and the explicit rate}\label{sec:removal}

\begin{proof}[Proof of Theorem~\ref{thm:removal}]
Let $\cA\subseteq[1,N]$ attain $F_V(N)$.
Delete its nonregular elements.
By Lemma~\ref{lem:exceptional}, this removes at most $500000z^{-K}N$ elements.
If a remaining reciprocal triple had primitive parameter $v\le V$, then
$u,v,u+v\le2V$.
All three are $P_V$-smooth, so the triple would already be forbidden in the relaxation.
Every remaining triple therefore has $v>V$ and is counted in Proposition~\ref{prop:regularcount}.

From each remaining triple delete one of its elements.
The union of all these chosen elements has cardinality at most the number of triples and meets every forbidden triple.
The resulting set is reciprocal-sum-free.
The total deletion cost is bounded by the right side of \eqref{eq:removal}.
The lower inequality follows from $f(N)\le F_V(N)$.
\end{proof}

\begin{proof}[Proof of the convergence assertion in Theorem~\ref{thm:density}]
Lemma~\ref{lem:fibreidentity} and $f(N)\le F_V(N)$ imply
\[
\limsup_{N\to\infty}\frac{f(N)}N\le U_V
\]
for every $V$, hence the upper limit is at most $C$.
On the other hand, Theorem~\ref{thm:removal} gives, with $V,K$ fixed,
\begin{equation}\label{eq:liminflower}
\liminf_{N\to\infty}\frac{f(N)}N
\ge U_V-500000z^{-K}-10^{23}q^{3K}(\log V)^{-\kappa}.
\end{equation}
Given $\eps>0$, first choose $K$ so that $500000z^{-K}<\eps/2$, then choose $V$ so that the second error in \eqref{eq:liminflower} is below $\eps/2$.
Since $U_V\ge C$, the lower limit is at least $C-\eps$.
Letting $\eps\downarrow0$ proves \eqref{eq:Cdef}.
\end{proof}

Passing to the limit in \eqref{eq:removal} now gives, for all $V\ge4$ and $K\ge0$,
\begin{equation}\label{eq:limitingremoval}
0\le U_V-C\le500000z^{-K}+10^{23}q^{3K}(\log V)^{-\kappa}.
\end{equation}

\begin{proof}[Proof of Corollary~\ref{cor:rates}]
The exact rational comparison $q^4<z^{15}$ implies
\begin{equation}\label{eq:logcomparison}
\frac{3\log q}{140\log z}<\frac9{112},
\qquad \frac7{80}-\frac9{112}=\frac1{140}.
\end{equation}
For $V\ge4$, choose
\[
K=\floor{\frac{\log\log V}{140\log z}}.
\]
Then
\[
z^{-K}\le z(\log V)^{-1/140},
\qquad
q^{3K}(\log V)^{-\kappa}\le(\log V)^{-1/140}.
\]
Insert these into \eqref{eq:limitingremoval} and use $500000z=561000$ to obtain \eqref{eq:outerrate}.
For \eqref{eq:finiterate}, make the same choice of $K$ with $V$ replaced by $\min(V,N)$.
Both logarithmic tails in \eqref{eq:removal} are then at most the tail at that minimum.
Division by $N$ completes the proof.
\end{proof}

\section{Effective determination of the constant}\label{sec:effective}

\begin{proposition}[Rational enclosures]\label{prop:enclosures}
For every rational $\eps>0$, finite integer and rational operations determine rational numbers $a_\eps\le C\le b_\eps$ with $b_\eps-a_\eps<\eps$.
\end{proposition}

\begin{proof}
It suffices to treat $0<\eps<1$.
Choose an integer $K\ge0$ such that $500000(500/561)^K\le\eps/4$.
Next choose an integer $B\ge1$ satisfying
\[
B^7\ge\left(\frac{4\cdot10^{23}q^{3K}}{\eps}\right)^{80},
\qquad V=2^{2B}.
\]
Both choices can be made by increasing an integer and testing rational inequalities.
Since $\log V>B$, \eqref{eq:limitingremoval} gives $U_V-C\le\eps/2$.

Enumerate the primes in $P_V$ and the $P_V$-smooth integers through an increasing bound $T$ until $E_{P_V}(T)\le\eps/4$, where $E_P$ is defined in \eqref{eq:prefix}.
This search terminates by \eqref{eq:prefixerror}.
For every smooth prefix, its independence number is determined by finite enumeration of subsets, testing the integer equality $bc=a(b+c)$ on triples of distinct elements.
Thus $L=L_{P_V}(T)$ is an explicitly determined rational number.

The inequalities $L\le U_V\le L+\eps/4$ and $U_V-\eps/2\le C\le U_V$ imply
\[
L-\eps/2\le C\le L+\eps/4.
\]
These endpoints have difference $3\eps/4<\eps$.
\end{proof}

\begin{proposition}[Finite-prime error]\label{prop:finiteerror}
If $P$ has $d$ primes and $U_P$ denotes the limit in \eqref{eq:fixedlimit}, then
\begin{equation}\label{eq:finiteerror}
\left|\frac{F_P(N)}N-U_P\right|\le\frac{8^d+4^d}{\sqrt N}\qquad(N\ge1).
\end{equation}
\end{proposition}

\begin{proof}
By geometric series,
\[
\sum_{s\in\Sp(P)}s^{-1/2}
=\prod_{p\in P}(1-p^{-1/2})^{-1}\le4^d.
\]
Consequently, $\#(\Sp(P)\cap[1,N])\le4^d\sqrt N$ and
$\sum_{s>N,\ s\in\Sp(P)}1/s\le4^d/\sqrt N$.
Apply \eqref{eq:roughcount} termwise to \eqref{eq:fibreidentity}.
The sum of its finite counting errors, after division by $N$, is at most $8^d/\sqrt N$.
The omitted tail is at most $4^d/\sqrt N$, since $\rho_P\le1$ and $\Delta_{P,j}\le1$.
\end{proof}

\begin{corollary}[An asymptotic modulus]\label{cor:modulus}
For every integer $N\ge\exp(40000)$,
\begin{equation}\label{eq:modulus}
\left|\frac{f(N)}N-C\right|
\le2(2\cdot10^{23}+561000)(\log\log N)^{-1/140}+2N^{-9/20}.
\end{equation}
\end{corollary}

\begin{proof}
Put $V=\floor{(\log N)/100}$ and $d=|P_V|\le2V$.
Since $\log8<5/2$, Proposition~\ref{prop:finiteerror} gives
\[
\left|\frac{F_V(N)}N-U_V\right|
\le2\exp\left(\frac{2\log8}{100}\log N-\frac12\log N\right)
\le2N^{-9/20}.
\]
Write $U=\log N\ge40000$.
Then $V\ge U/200$, so
$\log V\ge\log U-\log200\ge\tfrac12\log U$.
Also $4\le V\le N$.

Set $e=|F_V(N)/N-U_V|$, $a=U_V-C$, and $b=(F_V(N)-f(N))/N$.
The numbers $a,b$ are nonnegative, and
\[
\frac{f(N)}N-C=\left(\frac{F_V(N)}N-U_V\right)+a-b.
\]
Corollary~\ref{cor:rates} bounds both $a$ and $b$ by
$(2\cdot10^{23}+561000)(\log V)^{-1/140}$.
Hence the absolute value is at most $e$ plus that bound.
The preceding estimate for $\log V$, together with $2^{1/140}\le2$, proves \eqref{eq:modulus}.
\end{proof}

\section{The lower construction}\label{sec:lower}

The following construction is due to Stijn Cambie and is recorded with the original problem \cite{Bloom}.
We include its proof to settle the proposed value $1/2$ within the same framework.

\begin{proposition}\label{prop:cambie}
For every integer $N\ge1$,
\begin{equation}\label{eq:cambie}
f(N)\ge\floor{\frac{\floor{N/4}+1}{2}}+N-\floor{N/2}
\ge\frac58N-1.
\end{equation}
In particular, $C\ge5/8$ and $f(N)/N$ does not converge to $1/2$.
\end{proposition}

\begin{proof}
Let $\cA$ consist of the odd integers at most $N/4$, together with the integers in $(N/2,N]$.
These two parts are disjoint and have the cardinalities in \eqref{eq:cambie}.
The last inequality follows from the elementary bounds for integer parts.

Suppose $a,b,c\in\cA$ satisfy \eqref{eq:reciprocal}.
If $a>N/2$, then $1/a<2/N\le1/b+1/c$, a contradiction.
Thus $a$ is odd and $a\le N/4$.
If either $b$ or $c$ belongs to the lower part, it too is odd.
Reducing $bc=a(b+c)$ modulo $2$ then gives a contradiction.
The remaining possibility is $b,c>N/2$, in which case
$1/b+1/c<4/N\le1/a$, again impossible.
The set is reciprocal-sum-free.
Divide \eqref{eq:cambie} by $N$ and use Theorem~\ref{thm:density}.
\end{proof}

\subsection*{Acknowledgments}
The authors acknowledge substantial assistance from ChatGPT (OpenAI) in developing and checking the arguments and in preparing the exposition.

\begin{thebibliography}{9}

\bibitem{Bloom}
T. F. Bloom,
\emph{Erd\H{o}s problem 302},
\url{https://www.erdosproblems.com/302},
including the construction attributed to Stijn Cambie.

\bibitem{DP}
D. Della Pietra,
\emph{A positive-density improvement for all-length unit-fraction-free sets},
manuscript, July 30, 2026,
\href{https://github.com/donalddellapietra/erdos-301-proof/blob/789c6f045dbc81da3811031247d186a7128dafce/paper/main.tex}{version of July 30, 2026}.

\bibitem{EG}
P. Erd\H{o}s and R. L. Graham,
\emph{Old and New Problems and Results in Combinatorial Number Theory},
Monographies de L'Enseignement Math\'ematique, vol.~28,
Universit\'e de Gen\`eve, Geneva, 1980, p.~37.

\bibitem{CGM}
E. H. Jiang, J. Zhang, J. Yu, J. Lu, Z. Zhang, M. Simon, et al.,
\emph{Continual Graph Memory and Community Infrastructure for AI-Assisted Mathematical Discovery},
\href{https://arxiv.org/abs/2610.02945}{arXiv:2610.02945}, 2026.
See also the companion manuscript
\emph{Erd\H{o}s Problem \#302: A Consolidated Candidate Proof of the Density Limit},
\href{https://github.com/uclanlp/continual-graph-memory/blob/23367c6568c344d9c4dbf59be3f39b39f5d54324/results/erdos_302.md}{version of September 26, 2026},
Sections 2--6 and Appendix B.

\end{thebibliography}

\end{document}
